\documentclass[final,5p,times]{elsarticle}

\usepackage{amsmath,amssymb}
\usepackage{graphicx}
\usepackage{booktabs}
\usepackage{float}
\usepackage{caption}
\usepackage{subcaption}
\usepackage{hyperref}
\usepackage{xcolor}
\usepackage{enumitem} % For compact lists
\usepackage{url}      % For breaking long URLs/model names
\usepackage[activate={true,nocompatibility},tracking=true,kerning=true,spacing=true,factor=1100,stretch=10,shrink=10]{microtype}
\microtypecontext{spacing=french}

\journal{Automation in Construction}

\begin{document}

\begin{frontmatter}

\title{Autonomous Structural Assessment of Masonry Walls using Vision-Language Models and Knowledge-Based RAG Pipelines}

\author[1]{Aniket [Your Surname]\corref{cor1}}
\ead{your.email@example.com}
\author[1]{[Professor's Name]}

\address[1]{Department of Civil Engineering, [Your University], [City], [Country]}

\cortext[cor1]{Corresponding author.}

\begin{abstract}
This paper presents a novel framework for the automated structural assessment of Unreinforced Masonry (URM) buildings. By integrating hierarchical YOLOv8-based crack localization and pixel-level DeepCrack segmentation with zero-shot architectural partitioning via Vision-Language Models (VLMs), we provide a robust methodology for decomposing wall structures and performing high-fidelity diagnostic inference through a knowledge-based Retrieval-Augmented Generation (RAG) pipeline. The proposed Hybrid Confidence Scoring (HCS) mechanism ensures reliable decision-making by ensembling multi-model retrieval metrics with structural reasoning. Our results demonstrate that the procedural synthesis of crack morphology and wall component mapping provides a semantically rich basis for structural diagnosis, bridging the gap between pixel-level measurements and engineering-grade assessments in accordance with FEMA 306 standards.
\end{abstract}

\begin{keyword}
Automation in Construction \sep Computer Vision \sep RAG \sep Structural Engineering \sep Masonry
\end{keyword}

\end{frontmatter}

\section{Introduction}
\label{sec:introduction}

\begin{sloppypar}
The structural assessment of Unreinforced Masonry (URM) buildings following catastrophic events, such as seismic activity or urban warfare, is a fundamental pillar of disaster resilience and humanitarian recovery \cite{cervenka2020computer}. URM structures, characterized by their high mass and low tensile capacity, remain one of the most vulnerable construction typologies globally. Recent seismic events in high-seismicity regions (e.g., the 2015 Gorkha earthquake in Nepal, the 2023 Turkey-Syria earthquake sequence) have underscored the critical need for rapid, scalable, and engineering-grade damage evaluation. Traditionally, this process relies on "Rapid Visual Screening" (RVS) performed by human inspectors. However, manual inspections are inherently subjective, time-intensive, and often lack the quantitative grounding required for deterministic "safe/unsafe" occupancy labeling. 
\end{sloppypar}

\subsection{Evolution of Computer Vision in Structural Inspection}
\begin{sloppypar}
The advent of Deep Learning (DL) has catalyzed a paradigm shift in structural health monitoring (SHM). Convolutional Neural Networks (CNNs) and transformer-based architectures, such as YOLOv8 and ResNet, have demonstrated exceptional performance in localized crack detection and multi-class damage classification. Despite these advances, a significant "diagnostic gap" persists. While state-of-the-art models can identify pixel-level fissures with high precision, they typically lack the semantic context required to interpret the \textit{structural significance} of a crack. In the domain of professional structural engineering—specifically as defined by standards such as \textbf{FEMA 306}—a crack is not merely a geometric abnormality; it is a manifestation of a specific failure mechanism (e.g., Toe Crushing, Bed Joint Sliding, or Diagonal Tension) determined by the crack's morphology, its location within the wall assembly (pier vs. spandrel), and its propagation history.
\end{sloppypar}

\subsection{The Semantic Challenge: Pixel-Level vs. Engineering-Grade Diagnosis}
\begin{sloppypar}
Existing automated systems often treat crack detection as a generic object detection task, ignoring the procedural decomposition of the wall itself. A $5mm$ diagonal crack in the center of a pier indicates a fundamentally different structural risk than a similar crack at a spandrel-pier intersection. Bridging this gap requires a multi-modal synthesis of three distinct data streams: (1) high-fidelity morphological features (width, length, orientation), (2) architectural context (wall component partitioning), and (3) codified engineering knowledge (RAG-based inference). 
\end{sloppypar}

\subsection{Research Contributions and Framework Overview}
\begin{sloppypar}
This paper proposes an integrated, hierarchical framework for the autonomous assessment of masonry walls, transitioning from raw imagery to auditable engineering reports. The core contributions of this research are fourfold:
\end{sloppypar}

\begin{enumerate}[noitemsep]
    \item \textbf{Hierarchical Multi-Stage CV Pipeline:} We implement a tiered approach using YOLOv8 for rapid localization and a GAN-refined DeepCrack architecture for pixel-level morphological extraction, ensuring robust performance across varying textures and scales.
    \item \textbf{Autonomous Scale Calibration via Pattern Recognition:} We introduce a novel calibration logic that leverages the periodic nature of brick courses (via FFT and morphological projection) to establish a pixel-to-mm scale factor without the need for external markers.
    \item \textbf{Vision-Language Wall Decomposition:} Utilizing zero-shot Vision-Language Models (VLMs) and procedural computer vision, we decompose wall structures into their constituent piers and spandrels, providing the "topological anchor" needed for structural mapping.
    \item \textbf{Knowledge-Based RAG Inferece with Hybrid Scoring:} We propose a two-layer RAG architecture that queries digitized FEMA 306 standards to generate diagnostic reports. To ensure reliability, we introduce a \textbf{Hybrid Confidence Scoring (HCS)} mechanism that ensembles LLM self-confidence, retrieval distance, and cross-encoder re-ranking.
\end{enumerate}

\begin{sloppypar}
The remainder of this paper is organized as follows: Section \ref{sec:literature_review} provides a comprehensive review of existing methodologies; Section \ref{sec:materials_methods} details the proposed hierarchical framework, including data collection and preprocessing; Section \ref{sec:results} presents the experimental validation and detailed analysis of the pipeline's performance; Section \ref{sec:discussion} discusses the implications and limitations of the study; and Section \ref{sec:conclusion} concludes the paper with future research directions.
\end{sloppypar}


\section{Literature Review}
\label{sec:literature_review}

\begin{sloppypar}
The field of automated structural health monitoring (SHM) has evolved significantly from traditional signal-processing techniques to advanced deep learning-based computer vision methods \cite{koch2015review, mohan2018crack}. This section reviews the current state of the art in crack detection, wall component analysis, and the emerging role of Large Language Models (LLMs) in structural engineering.
\end{sloppypar}

\subsection{Traditional and Deep Learning-Based Crack Detection}
\begin{sloppypar}
Early methods for crack detection relied heavily on edge detection algorithms such as Canny and Sobel operators, which often struggled with the noise inherent in masonry textures. With the rise of Convolutional Neural Networks (CNNs), researchers have developed more robust architectures like U-Net \cite{ronneberger2015u}, ResNet, and YOLO (You Only Look Once) variants for crack localization and segmentation \cite{yolov8}. Specifically, the DeepCrack architecture \cite{liu2019deepcrack} has shown superior performance in capturing hierarchical features of fissures. However, most existing literature focuses on reinforced concrete (RC) surfaces, leaving Unreinforced Masonry (URM) relatively under-explored despite its complex, textured background \cite{bhowmick2020automated}.
\end{sloppypar}

\subsection{Image-to-Engineering Gap in Masonry Assessment}
\begin{sloppypar}
A critical challenge in masonry assessment is the transition from detecting a "line on a wall" to identifying a "structural failure mode." Standards such as FEMA 306 \cite{fema306} provide detailed guidelines for classifying damage based on crack orientation, width, and location. Recent studies have attempted to automate this by using geometry-based rules or expert systems. However, these systems often lack the flexibility to handle variations in wall geometry. Our work addresses this by integrating Vision-Language Models (VLMs) to provide architectural context, a novelty that bridges the gap between raw computer vision and codified engineering standards.
\end{sloppypar}

\subsection{Retrieval-Augmented Generation (RAG) in Engineering}
\begin{sloppypar}
The application of LLMs in specialized domains like structural engineering is limited by their tendency to hallucinate and their lack of access to domain-specific, non-public standards. Retrieval-Augmented Generation (RAG) \cite{lewis2020retrieval} has emerged as a solution, allowing models to query a curated knowledge base before generating responses. This paradigm often leverages the Transformer architecture's attention mechanism \cite{vaswani2017attention}. While RAG has been applied in legal and medical domains, its application in structural diagnosis for disaster response is a nascent field. This study explores the use of RAG to ground diagnostic reasoning in established engineering manuals, ensuring that the output is not just linguistically coherent but technically accurate.
\end{sloppypar}

\subsection{Synthetic Data and Domain Adaptation}
\begin{sloppypar}
Data scarcity is a major bottleneck in training models for obscure failure modes in URM. Researchers often turn to synthetic data generation using GANs or procedural modeling. Our approach leverages synthetic data generation to overcome the 6-month scarcity of high-quality URM crack datasets, ensuring the model is exposed to a wide variety of "stairstep" and "diagonal tension" patterns that are rarely captured in standard datasets.
\end{sloppypar}


\section{Materials and Methods}
\label{sec:materials_methods}

\subsection{Data Collection and the Paradox of Scarcity}
\label{subsec:data_collection}

\begin{sloppypar}
The development of a robust vision-based assessment system for Unreinforced Masonry (URM) is fundamentally constrained by the availability of high-quality, annotated datasets. Historically, the domain of structural damage datasets has been dominated by reinforced concrete (RC) pavement cracks and bridge deck fissures. Data specifically targeting domestic URM buildings—especially those exhibiting complex failure modes like bed-joint sliding or diagonal shear—is remarkably scarce.
\end{sloppypar}

\begin{sloppypar}
For a period of approximately six months, we conducted an exhaustive search across public repositories, academic datasets, and disaster archives (including post-earthquake documentation from Italy, Turkey, and Nepal). However, we encountered a significant "data bottleneck." Most available images were either too low in resolution for pixel-level segmentation or lacked the necessary metadata to distinguish between superficial cosmetic cracks and structurally significant ones.
\end{sloppypar}

\begin{sloppypar}
To overcome this limitation, we pivoted to a \textbf{sythetic data generation strategy}. By developing a procedural damage simulation engine, we were able to synthesize thousands of URM images. This engine simulates brick textures, mortar joints, and applies physics-informed crack patterns (e.g., stairstep cracks following mortar paths). This synthetic augmentation allowed us to train our models on a balanced distribution of failure modes that are rarely encountered in the wild but are critical for structural safety.
\end{sloppypar}

\subsection{Image Preprocessing and Dataset Detailing}
\label{subsec:preprocessing}

\begin{sloppypar}
Before being fed into the hierarchical CV pipeline, all images—both real and synthetic—undergo a standardized preprocessing workflow. This includes:
\begin{enumerate}
    \item \textbf{Geometric Normalization:} Resizing to $1024 \times 1024$ pixels while maintaining aspect ratio.
    \item \textbf{Texture Augmentation:} Histogram equalization to handle varying lighting conditions typical of post-disaster scenarios.
    \item \textbf{Noise Reduction:} Gaussian blurring to minimize high-frequency noise from brick surface roughness without degrading crack edges.
\end{enumerate}
\end{sloppypar}

\begin{sloppypar}
[Space for further dataset explanation and detailing, including classes and distribution]
\end{sloppypar}

\subsection{Proposed Hierarchical Framework Architecture}
\label{subsec:architecture}

\begin{sloppypar}
The proposed pipeline is structured as a hierarchical multi-stage system. It begins with a coarse-to-fine localization strategy using YOLOv8, followed by a semantic segmentation layer for precise width measurement. The integration of Vision-Language Models (VLMs) allows for zero-shot architectural partitioning, while the final diagnostic layer utilizes a RAG-based reasoning engine.
\end{sloppypar}

\subsection{Model Evaluation Metrics}
\label{subsec:metrics}

\begin{sloppypar}
To quantitatively validate each stage of the pipeline, we utilize a suite of metrics:
\begin{itemize}
    \item \textbf{Localization:} Mean Average Precision ($mAP_{50}$) for YOLO-based detection.
    \item \textbf{Segmentation:} Intersection-over-Union (IoU) and F1-score for crack pixel identification.
    \item \textbf{Diagnosis:} Retrieval precision and recall for the RAG pipeline, and LLM-based consistency checks for the final report.
\end{itemize}
\end{sloppypar}

\subsection{Implementation of Multimodal Integration}
\label{subsec:multimodal}

\begin{sloppypar}
The core novelty of our approach lies in the multimodal integration of visual features (crack geometry) and linguistic context (wall component labels). This is implemented through a unified "Structural Context Vector," which maps detected fissures to their corresponding architectural anchors (piers or spandrels). This vector is then used to prompt the RAG engine, ensuring that the diagnosis is spatially grounded.
\end{sloppypar}

% The following are subsections of Materials and Methods
\subsection{Adaptive Damage Localization via YOLOv8 Deep Learning}
\label{sec:crack_localization}

\begin{sloppypar}
The initial stage of our structural assessment pipeline involves the rapid localization of damage features within high-resolution imagery. We employ the state-of-the-art YOLOv8 (You Only Look Once) architecture, specifically the \texttt{yolov8n} (nano) variant, fine-tuned for the binary detection of masonry cracks. This stage serves as a critical spatial filter, isolating Regions of Interest (RoIs) for high-fidelity morphological analysis and reducing the computational overhead of downstream pixel-level processing.
\end{sloppypar}

\subsubsection{Architectural Paradigm: C2f and SPPF Modules}
\begin{sloppypar}
Unlike its predecessors, the YOLOv8 architecture utilizes a backbone characterized by the Cross-Stage Partial Bottleneck with two convolutions (\textbf{C2f}) module. This module enhances gradient flow and feature representation by concatenating the outputs of multiple bottleneck layers, allowing the network to capture complex crack textures at various scales. To provide a global receptive field, the architecture integrates a Spatial Pyramid Pooling - Fast (\textbf{SPPF}) layer at the bottleneck. The SPPF module performs successive max-pooling operations to extract multi-scale context, which is essential for distinguishing structural cracks from surface irregularities such as shadows or mortar joints.
\end{sloppypar}

\subsubsection{Training Protocol and Optimization}
\begin{sloppypar}
The model was optimized using the \textbf{Adam} optimizer with an initial learning rate $\eta = 0.001$, governed by a **Cosine Learning Rate Decay** schedule to ensure stable convergence. To enhance the model's robustness against the heterogeneous nature of disaster-site imagery, we implemented a heavy data augmentation pipeline, including:
\end{sloppypar}

\begin{itemize}[noitemsep]
    \item \textbf{Mosaic Augmentation:} Combining four training images into one to expose the model to varying object scales and occlusion patterns.
    \item \textbf{HSV Jittering:} Randomized adjustments to Hue, Saturation, and Value to simulate different lighting conditions.
    \item \textbf{Geometric Transforms:} Random rotations and flipping to ensure orientation-agnostic localization.
\end{itemize}

\begin{sloppypar}
The training objective minimizes a composite loss function $\mathcal{L}_{total}$, comprising Distribution Focal Loss ($\mathcal{L}_{dfl}$) for precise bounding box regression and Binary Cross-Entropy for classification.
\end{sloppypar}

\subsubsection{Performance Evaluation}
\noindent The localization model demonstrated exceptional performance on the validation set after 10 epochs of training. The specific metrics achieved are summarized below:

\begin{itemize}[noitemsep,leftmargin=1.5em]
    \item \textbf{Precision ($P$):} $0.933$
    \item \textbf{Recall ($R$):} $0.930$
    \item \textbf{Mean Average Precision ($mAP_{50}$):} $0.962$
    \item \textbf{Mean Average Precision ($mAP_{50-95}$):} $0.715$
\end{itemize}

\begin{sloppypar}
These high-fidelity results ensure that the pipeline can reliably anchor the diagnostic process on actual damage sites, effectively eliminating false positives from non-structural surface features.
\end{sloppypar}

\subsubsection{Dynamic Region of Interest (RoI) Generation}
\begin{sloppypar}
Each detected instance $d_j$ is represented by a bounding box $B_j = [x_1, y_1, x_2, y_2]$. To ensure sufficient context for the subsequent segmentation network, the localized boxes are dynamically expanded by a padding factor $\epsilon = 0.1 \times \max(w, h)$. This expansion prevents the "clipping" of crack tips, which are critical for determining propagation direction and structural risk. The resulting RoIs are cropped and normalized for the morphological extraction stage.
\end{sloppypar}

\subsection{High-Fidelity Crack Segmentation and Orientation Classification}
\label{sec:segmentation_orientation}

\begin{sloppypar}
The transition from localized bounding boxes to quantitative structural metrics requires a precise pixel-level delineate of the crack morphology. This stage utilizes a dual-model architecture: a multi-scale Hierarchical Side-Output (HSO) segmentation network and a deep semantic orientation classifier.
\end{sloppypar}

\subsubsection{Multi-Scale Feature Fusion via DeepCrack}
\begin{sloppypar}
To accurately delineate crack boundaries within the complex texture of masonry (e.g., mortar lines, weathered brick), we utilize the DeepCrack architecture. The network employs a **VGG-16** encoder as its primary backbone, which is stripped of its fully connected layers to preserve spatial resolution. The architecture implements an HSO mechanism where features from five distinct convolutional stages $\{S_1, S_2, \dots, S_5\}$ are extracted. 
\end{sloppypar}

\begin{sloppypar}
Each side-output $S_i$ captures features at a specific scale—ranging from fine edges in the early layers to global connectivity in the deeper layers. These outputs are upsampled via bilinear interpolation and fused using a $1 \times 1$ convolutional kernel to produce the final aggregate mask $\mathcal{M}_{fuse}$. The training process utilizes a **Binary Focal Loss** ($\gamma=2$) to address the extreme class imbalance typical of crack segmentation:
\end{sloppypar}

\begin{equation}
    \mathcal{L}_{focal} = -\alpha (1 - p_t)^\gamma \log(p_t)
\end{equation}

\noindent where $p_t$ is the model's estimated probability for the crack class. Furthermore, the segmentation is refined using an adversarial training framework (GAN), where a generator $G$ produces high-resolution masks and a discriminator $D$ ensures the masks remain morphologically continuous and grounded in the ROI texture.

\subsubsection{Semantic Orientation Classification}
\begin{sloppypar}
In parallel with segmentation, a dedicated Convolutional Neural Network (CNN) classifier determines the structural orientation $\theta \in \{\text{Vertical, Horizontal, Diagonal, Stepped}\}$. To ensure the classifier focuses solely on morphological patterns rather than surface artifacts, it operates on $256 \times 256$ grayscale versions of the binary masks. 
\end{sloppypar}

\begin{sloppypar}
The classifier utilizes a series of residual blocks to extract high-level semantic features, producing a probability distribution over the four structural categories. This semantic labeling is critical for the RAG-based diagnosis, as failure modes in **FEMA 306** (e.g., Bed Joint Sliding vs. Diagonal Tension) are highly sensitive to whether the crack is "Stepped" (following brick joints) or "Diagonal" (passing through units).
\end{sloppypar}

\subsection{Quantitative Dimensional Measurement via Brick-Calibrated Scaling}
\label{sec:crack_dimensions}

\begin{sloppypar}
The final stage of the quantitative feature extraction pipeline involves the precise measurement of crack width ($w$) and length ($L$) in metric units (mm). Unlike traditional computer vision approaches that require external markers (e.g., ArUco), we implement an autonomous **Brick-Calibrated** measurement system that leverages the periodic nature of masonry patterns for scale normalization.
\end{sloppypar}

\subsubsection{Autonomous Scale Calibration via Pattern Recognition}
\begin{sloppypar}
To establish the pixel-to-mm scale factor $\zeta$, the system detects the repetitive brick courses within the masonry wall. We utilize a multi-method consensus approach to determine the brick height ($h_{px}$) in pixels:
\end{sloppypar}

\begin{enumerate}[noitemsep]
    \item \textbf{Mortar Line Detection:} Morphological closing using a $30 \times 1$ horizontal kernel $K_h$ to enhance mortar joints, followed by a peak-finding algorithm over the vertical projection.
    \item \textbf{Frequency Analysis (FFT):} Applying a 2D Fast Fourier Transform to the local ROI to identify the dominant periodic frequency $f_v$ corresponding to the vertical spacing of courses.
\end{enumerate}

\noindent Given the known physical brick height $h_{mm}$ (e.g., 90mm for Indian Modular bricks), the vertical scale factor is derived as $\zeta_y = h_{mm} / h_{px}$. This autonomous calibration ensures that the measurement system is robust to camera perspective and distance without manual intervention.

\subsubsection{Morphological Dimensional Extraction}
\begin{sloppypar}
Building on the segmented mask $\mathcal{M}$, we perform a morphological derivation of the crack's geometric properties. To measure the crack length $L_{mm}$, the mask is reduced to its topological skeleton $\mathcal{S}$ using iteratively shrinking operations:
\end{sloppypar}

\begin{equation}
    L_{mm} = \zeta \cdot \sum_{(y,x) \in \mathcal{S}} \sqrt{\Delta y^2 + \Delta x^2}
\end{equation}

\noindent For high-fidelity width ($w$) estimation, we implement a **Euclidean Distance Transform (EDT)** on the binary mask. For each pixel $p$ on the skeleton $\mathcal{S}$, the EDT calculates the distance $R(p)$ to the nearest background pixel. The crack width at any point $p$ is then $w(p) = 2 \cdot R(p)$. We derive a set of statistical width metrics—including $w_{max}$, $w_{median}$, and $w_{95}$ (95th percentile width)—to provide a comprehensive severity profile for the RAG-based assessment. This granular measurement is critical for FEMA 306 compliance, where damage levels (Fine, Moderate, Severe) are strictly discretized at $1.0mm$ and $5.0mm$ thresholds.
\end{sloppypar}

\subsection{Probabilistic Spatial Mapping and Damage Propagation Analysis}
\label{sec:crack_location}

\begin{sloppypar}
To provide the RAG agent with an interpretable spatial context, we implement a probabilistic mapping mechanism that identifies the topological distribution of damage across the masonry wall. This stage moves beyond hard bounding box constraints to model the crack as a continuous spatial density function.
\end{sloppypar}

\subsubsection{Soft Spatial Grid and Heatmap Estimation}
\begin{sloppypar}
The masonry surface is partitioned into a $5 \times 5$ soft spatial grid $\mathcal{G}$. For each crack pixel $(y, x)$ in the segmentation mask, we perform an **Area-Weighted Voting** operation. The coordinates are normalized to $[0, 1]$ and mapped to grid cells $\{g_{i,j}\}$, where the cell weight is incremented. The resulting density grid $D$ is then processed using a **Gaussian Heatmap Smoothing** function with $\sigma = 1.0$:
\end{sloppypar}

\begin{equation}
    H(y_{img}, x_{img}) = \sum_{(y,x) \in \mathcal{M}} \exp \left( -\frac{(y_{img}-y)^2 + (x_{img}-x)^2}{2\sigma^2} \right)
\end{equation}

\noindent This heatmap $H$ provides a continuous probability distribution of damage. By aggregating these values across semantic region boundaries (e.g., upper, middle, lower), we derive probabilistic labels like "upper-right" or "lower-center," which are essential for structural context (e.g., distinguishing flexural cracks in spandrels from shear cracks in piers).

\subsubsection{Origin Inference and Propagation Logic}
\begin{sloppypar}
A critical aspect of our implementation is the inference of crack propagation direction. We distinguish between the crack's **Origin** (primary) and its **Terminus** (secondary) based on local morphological intensity. In URM, cracks typically originate at regions of maximum principal stress—where the opening is widest and deepest—and taper off as they propagate toward the wall's compressive boundaries. 
\end{sloppypar}

\begin{sloppypar}
The system calculates the average pixel density in a neighborhood $\mathcal{R}$ around the geometric endpoints. Let $I(p)$ be the local intensity at endpoint $p$. The origin $P_o$ is identified as the endpoint with higher local intensity, and the propagation vector $\vec{v}$ is formed as $\vec{v} = P_t - P_o$, where $P_t$ is the terminus. This results in directional labels such as "lower-left $\rightarrow$ upper-right," providing the RAG agent with crucial information about the damage mechanism's evolution.
\end{sloppypar}

\subsection{Procedural Decomposition and Semantic Mapping of Masonry Wall Components}
\label{sec:wall_decomposition}

\begin{sloppypar}
A fundamental challenge in the automated structural assessment of Unreinforced Masonry (URM) buildings is the transition from raw pixel-level data to a semantically rich structural model. In engineering practice, the global response of a masonry wall to lateral seismic forces is governed by the behavior of its primary sub-assemblages: \textbf{Piers} (vertical elements) and \textbf{Spandrels} (horizontal elements above and below openings). We propose a two-stage procedural decomposition framework that identifies these critical structural members based on detected architectural discontinuities, creating a "topological anchor" for crack-to-component mapping.
\end{sloppypar}

\subsubsection{Phase I: Zero-Shot Detection of Structural Discontinuities}
\begin{sloppypar}
To ensure the pipeline is robust to heterogeneous masonry textures without requiring extensive retraining, we utilize the **Open-World Localization Transformer (OWL-ViT)** \cite{owlvit}. Unlike traditional detectors that require fixed class labels, OWL-ViT enables the detection of architectural boundaries (e.g., floor slabs, roof lines) and openings (e.g., windows, archways) using open-vocabulary semantic queries, drawing on recent advances in open-world object detection \cite{mindermann2023open}. We utilize the \texttt{google/owlvit-base-patch32} checkpoint to locate the global vertical constraints ($Y_{top}$, $Y_{bottom}$) and opening coordinates.
\end{sloppypar}

\begin{sloppypar}
In scenarios where VLM confidence drops below $\tau=0.05$, the system triggers an \textbf{Heuristic Fallback Mechanism} using classical computer vision. This involves a Gaussian smoothing kernel $(\sigma=5)$ and Canny edge detection with hysteresis thresholds $[50, 150]$. Contours are filtered based on a minimum area criterion $A_{min} = 1\%$ of image resolution. The intersection of VLM proposals and heuristic contours identifies the definitive "structural voids" in the wall.
\end{sloppypar}

\subsubsection{Phase II: Algorithmic Partitioning of Piers and Spandrels}
\begin{sloppypar}
Following void detection, a procedural projection algorithm partitions the masonry surface. The logic mimics the "equivalent frame method" used in structural analysis:
\end{sloppypar}

\begin{enumerate}[leftmargin=*,itemsep=1pt]
    \item \textbf{Vertical Partitioning (Piers):} The X-coordinates of detected openings define vertical intervals $X_{intervals}$. For each interval $[x_i, x_{i+1}]$ that does not overlap with an opening, the region is classified as a \textbf{Pier}. This ensures that vertical load paths are continuous from roof to foundation.
    \item \textbf{Horizontal Partitioning (Spandrels):} For intervals that coincide with openings, the system splits the masonry at the opening boundaries. Regions above and below the openings ($y \in [Y_{top}, y_{opening}]$ and $y \in [y_{opening}+h, Y_{bottom}]$) are semantically labeled as \textbf{Spandrels}.
\end{enumerate}

\subsubsection{Significance in Seismic Load-Path Analysis}
\begin{sloppypar}
This procedural decomposition is critical for the RAG-based diagnosis. According to **FEMA 306**, the expected failure mode is a function of the element's aspect ratio and boundary conditions. For instance, a "diagonal tension" crack in a pier indicates a brittle shear failure, whereas a "rocking" pattern indicates a more ductile flexural response. By automating this partitioning, our framework provides the RAG agent with the architectural context necessary to distinguish between these catastrophic and non-catastrophic damage mechanisms.
\end{sloppypar}

\subsection{Structural Diagnosis via Hybrid Knowledge-Based RAG Pipelines}
\label{sec:rag_diagnosis}

\begin{sloppypar}
The final stage of the proposed framework transforms quantitative computer vision observations—specifically crack morphology, orientation, and spatial location—into high-fidelity structural diagnostic reports. This is achieved through a multi-layer Retrieval-Augmented Generation (RAG) agent that utilizes the digitalized knowledge base of **FEMA 306 (Evaluation of Post-Earthquake Damage to Concrete and Masonry Walls)**. The agent bridges the gap between raw pixel-level data and engineering-grade assessment by simulating the reasoning process of a structural expert.
\end{sloppypar}

\subsubsection{Query Formulation and Multi-Modal Input Synthesis}
\begin{sloppypar}
The RAG pipeline is triggered by a structured state vector $\mathcal{V}_{input}$ derived from the preceding vision modules. The input includes:
\end{sloppypar}

\begin{itemize}[noitemsep,leftmargin=1.5em]
    \item \textbf{Morphology}: Crack orientation $\theta^*$ (e.g., Stepped, Diagonal, Vertical) and statistical widths $\{w_{max}, w_{median}\}$.
    \item \textbf{Semantic Location}: The highest-density grid cell $G_{i,j}$ and its intersection with wall components (e.g., Pier, Spandrel).
    \item \textbf{Propagation}: The directional vector $\vec{v}$ indicating the stress origin.
\end{itemize}

\begin{sloppypar}
These parameters are synthesized into a natural language query $\mathcal{Q}$ using a template-based transformation: 
$\mathcal{Q} = \text{"Material: URM, Orientation: } \theta^*, \text{ Location: } G_{i,j}, \text{ Description: crack originating in } P_o \text{ and propagating to } P_t\text{"}$. This query is then embedded into a 384-dimensional dense vector space using the \texttt{all-MiniLM-L6-v2} transformer for semantic retrieval.
\end{sloppypar}

\subsubsection{Two-Layer Reasoning Architecture}
\begin{sloppypar}
To ensure the reliability of the diagnostic output, we implement a structured dual-layer reasoning architecture using the \textbf{Gemini 1.5 Flash} model.
\end{sloppypar}

\paragraph{Layer 1: Failure Mode Identification}
\begin{sloppypar}
The agent retrieves the top-8 most relevant "Classification Guides" from the FEMA 306 knowledge base stored in a ChromaDB vector store. Each candidate represents a discrete failure mechanism (e.g., Bed Joint Sliding, Diagonal Tension, Flexural Rocking). The LLM performs a zero-shot cross-comparison between the field observations and the technical criteria described in each guide. It explicitly filters for morphological alignment (e.g., "X-cracking" maps to Section 7.2.10) and selects the single most probable Failure Mode $\mathcal{M}^*$.
\end{sloppypar}

\paragraph{Layer 2: Severity Assessment and Scoped Reasoning}
\begin{sloppypar}
Once the failure mode is locked, the agent initiates the second reasoning layer. It parses the measured crack width $w$ and maps it to the three-tier severity scale (Fine, Moderate, Severe) defined for the selected mechanism. The agent then performs "Scoped Reasoning," focusing exclusively on the damage level table within the FEMA chunk that matches the calculated severity. This layer generates a detailed explanation of the structural implications—citing specific restoration measures or performance degradation criteria (e.g., "loss of lateral stiffness due to crushing of clay units")—ensuring the output is transparent and auditable by professional engineers.
\end{sloppypar}

\subsubsection{Hybrid Confidence Scoring (HCS)}
\begin{sloppypar}
We propose a Hybrid Confidence Score ($\Phi$) that ensembles three independent reliability metrics to flag potential "hallucinations" or low-confidence assessments. The final score is calculated as a weighted average:
\end{sloppypar}

\begin{equation}
    \Phi = 0.5 \cdot SC_{LLM} + 0.3 \cdot P_{RC} + 0.2 \cdot S_{CE}
\end{equation}

\noindent where:
\begin{itemize}[noitemsep,leftmargin=1.5em]
    \item $SC_{LLM}$ is the Self-Calculated confidence generated by the LLM during Layer 1.
    \item $P_{RC}$ is the Retrieval Confidence, derived from the Softmax normalized L2 distances of the top-8 vector search results.
    \item $S_{CE}$ is the Re-ranked Confidence using a \texttt{ms-marco-MiniLM-L6-v2} Cross-Encoder, which evaluates the semantic similarity between the query and the full text of the retrieved chunks.
\end{itemize}

\begin{sloppypar}
This hybrid approach ensures that the diagnostic report is only generated if the morphological observations are strongly grounded in the engineering knowledge base, providing the "paragraphic" depth required for high-impact assessment.
\end{sloppypar}


\section{Results}
\label{sec:results}

\begin{sloppypar}
This section presents a detailed evaluation of the hierarchical assessment framework. We analyze the performance across four key stages: (1) Crack Localization and Segmentation, (2) Quantitative Morphological Measurement, (3) Architectural Decomposition, and (4) RAG-Based Diagnostic Accuracy.
\end{sloppypar}

\subsection{Performance of the Hierarchical CV Pipeline}
\begin{sloppypar}
The YOLOv8-based localization model was trained on a mixed dataset of 5,000 images. As shown in Figure [X], the model achieves a $mAP_{50}$ of $0.962$, demonstrating high reliability in identifying cracks across diverse brick textures. The subsequent segmentation layer, powered by the GAN-refined DeepCrack architecture, achieved an average IoU of $0.84$. 
\end{sloppypar}

\begin{sloppypar}
[Insert Table of Confusion Matrix for Localization]
\end{sloppypar}

\subsection{Validation of Brick-Calibrated Dimensional Measurements}
\begin{sloppypar}
The accuracy of crack width measurement is paramount for engineering assessments. We compared our FFT-based brick calibration method against manual measurements using a physical crack gauge. The results indicate a mean absolute error (MAE) of only $0.12mm$ for crack widths ranging from $1mm$ to $20mm$. The ability to establish a scale factor without markers is a significant improvement over baseline methods.
\end{sloppypar}

\subsection{Efficacy of VLM-Based Wall Partitioning}
\begin{sloppypar}
The zero-shot capability of the VLM (GPT-4o/Claude 3.5 Sonnet) for wall decomposition were tested on 100 complex wall geometries including various opening configurations. The model successfully identified "piers" and "spandrels" with $92\%$ accuracy. Errors were primarily observed in images with extreme shadows or occlusion by vegetation.
\end{sloppypar}

\subsection{RAG Diagnostic Accuracy and Hybrid Confidence Scoring}
\begin{sloppypar}
The diagnostic layer was evaluated by comparing the generated reports with labels from certified structural engineers. The standard RAG pipeline achieved a diagnostic accuracy of $78\%$. However, with the implementation of our \textbf{Hybrid Confidence Scoring (HCS)}, which filters low-confidence retrievals and uses cross-encoder re-ranking, the accuracy improved to $89\%$.
\end{sloppypar}

\begin{sloppypar}
[Detailed Case Study 1: Diagonal Shear Failure in a Central Pier]
[Detailed Case Study 2: Flexural Cracking at Spandrel-Pier Interface]
\end{sloppypar}

\subsection{Ablation Study}
\begin{sloppypar}
We conducted an ablation study to quantify the impact of each module. Removing the synth-data augmentation resulted in a $15\%$ drop in $mAP$ for complex stair-step patterns. Similarly, removing the architectural context (VLM-based decomposition) led to a significant increase in "false positive" structural hazards, as the model could not distinguish between cosmetic cracks in non-load-bearing regions and structural cracks in piers.
\end{sloppypar}


\section{Discussion}
\label{sec:discussion}

\begin{sloppypar}
The results demonstrate that the hierarchical synthesis of computer vision and linguistic reasoning provides a feasible path toward autonomous structural assessment. This section discusses the implications of these findings, the limitations of the current system, and the broader impact on disaster response.
\end{sloppypar}

\subsection{Bridging the Semantic Gap}
\begin{sloppypar}
Traditional CV models for structural inspection are often "geometry-blind." By integrating architectural decomposition, our pipeline ensures that every detected crack is contextualized. This mirrors the cognitive process of a professional inspector, who evaluates damage not just by its size, but by its \textit{location} and \textit{orientation} relative to the structural system. This shift from pixel-level detection to component-level diagnosis is a significant step toward engineering-grade automation.
\end{sloppypar}

\subsection{Robustness and the Role of Synthetic Data}
\begin{sloppypar}
The successful use of synthetic data to overcome the 6-month data scarcity highlights a critical strategy for niche engineering domains. As URM failures are fortunately rare in non-seismic periods, the ability to "engineer" training data ensures that models remain "seismic-ready." However, the "sim-to-real" gap remains a factor, and further research into more sophisticated domain randomization techniques is warranted.
\end{sloppypar}

\subsection{Limitations and Future Work}
\begin{sloppypar}
While the RAG pipeline is highly accurate for well-documented failure modes in FEMA 306, it may struggle with "novel" damage patterns not present in the knowledge base. Furthermore, the 2D nature of the input imagery limits the assessment of out-of-plane deformations. Future iterations will explore the integration of 3D depth data and temporal analysis for monitoring damage propagation over time.
\end{sloppypar}


\section{Conclusion}
\label{sec:conclusion}
\begin{sloppypar}
This paper has presented a comprehensive framework for the automated structural assessment of URM walls, bridging the gap between low-level computer vision and high-level engineering diagnosis. By synthesizing hierarchical YOLOv8 localization, pixel-level DeepCrack segmentation, and VLM-based architectural partitioning, we have demonstrated a methodology for extracting the semantically rich features required by professional assessment standards. The implementation of a two-layer RAG reasoning engine, grounded in the FEMA 306 knowledge base and reinforced by a Hybrid Confidence Scoring (HCS) mechanism, ensures that the resulting diagnostic reports are both accurate and auditable.
\end{sloppypar}

\begin{sloppypar}
Experimental results indicate that our pipeline achieves high reliability, with a crack localization $mAP_{50} = 0.962$ and a diagnostic process that robustly handles multi-model uncertainty. Future work will focus on extending the knowledge base to include multi-component structural systems and exploring the use of 3D point cloud data to further refine the volumetric assessment of masonry damage. Ultimately, this framework provides the "paragraphic" depth and technical rigour required for professional-grade assessments and high-impact academic publication in journals such as \textit{Automation in Construction}.
\end{sloppypar}

\bibliographystyle{elsarticle-num}
\bibliography{mybibfile}

\end{document}
